\PassOptionsToPackage{dvipsnames,svgnames,table}{xcolor}
\documentclass[11pt]{article}
\usepackage[a4paper,margin=22mm]{geometry}
\usepackage[no-math]{fontspec}
\setmainfont{DejaVu Serif}
\usepackage{amsmath,amssymb,amsthm,mathtools,graphicx,xcolor,hyperref,xurl,xspace,ifthen,enumerate}
\usepackage[shortlabels]{enumitem}
\setlength{\emergencystretch}{4em}
\AtBeginDocument{\special{pdf:mapline pzdr Dingbats <uzdr.pfb}}
\SetMathAlphabet{\mathbf}{normal}{OT1}{cmr}{bx}{n}
\SetMathAlphabet{\mathbf}{bold}{OT1}{cmr}{bx}{n}
\newtheorem{theo}{Theorem}[section]
\newtheorem{lemm}[theo]{Lemma}
\newtheorem{thm}[theo]{Theorem}
\newtheorem{proposition}[theo]{Proposition}
\newtheorem{corollary}[theo]{Corollary}
\newtheorem{definition}[theo]{Definition}
\newtheorem{theorem}[theo]{Theorem}
\newtheorem*{remas*}{Remarks}
\newtheorem{remas}[theo]{Remarks}
\newtheorem*{exam*}{Example}
\newtheorem{ques}[theo]{Question}
\providecommand{\og}{«\nobreakspace}
\providecommand{\fg}{\nobreakspace»}
\newtheorem{lem}[theo]{Lemma}
\newtheorem{lemma}[theo]{Lemma}
\newtheorem{prop}[theo]{Proposition}
\newtheorem{coro}[theo]{Corollary}
\newtheorem{corol}[theo]{Corollary}
\newtheorem*{conj*}{Conjecture}
\newtheorem{conj}[theo]{Conjecture}
\newtheorem{conjecture}[theo]{Conjecture}
\newtheorem{Proposition}[theo]{Proposition}
\newtheorem{theoreme}[theo]{Theorem}
\theoremstyle{definition}
\newtheorem{defi}[theo]{Definition}
\newtheorem{exam}[theo]{Example}
\newtheorem{exem}[theo]{Example}
\newtheorem{remark}[theo]{Remark}
\newtheorem{example}[theo]{Example}
\newtheorem{exams}[theo]{Examples}
\newtheorem{rema}[theo]{Remark}
\newtheorem{nota}[theo]{Notation}
\newtheorem{notation}[theo]{Notation}
\newtheorem*{notas}{Notations}
\newtheorem{result}[theo]{Result}
\newtheorem{question}[theo]{Question}
\newtheorem{prob}[theo]{Problem}
\newtheorem*{rema*}{Remark}
\newtheorem*{defi*}{Definition}
\newtheorem*{theo*}{Theorem}
\newtheorem*{lemm*}{Lemma}
\newtheorem*{prop*}{Proposition}
\newcommand{\enoncename}{}
\newtheorem{corpusenonce}[theo]{\enoncename}
\newenvironment{enonce}[2][]{\renewcommand{\enoncename}{#2}\begin{corpusenonce}}{\end{corpusenonce}}
\newenvironment{enonce*}[2][]{\par\medskip\noindent\textbf{#2.}\itshape}{\par\medskip}
\newenvironment{proofc}[1][\proofname]{\begin{proof}[#1]}{\end{proof}}
\newcommand{\Mc}{\mathcal{M}}
\newcommand{\pointir}{.\ ---\ }
\newcommand{\equalenv}[2]{\expandafter\let\csname#1\expandafter\endcsname\csname#2\endcsname\expandafter\let\csname end#1\expandafter\endcsname\csname end#2\endcsname}
\providecommand{\diagbox}[3][]{\begin{tikzpicture}[x=1em,y=1em,baseline=(current bounding box.center)]\useasboundingbox(0,0) rectangle(3,1.4);\draw(0,1.4)--(3,0);\node at(.5,.3){#2};\node at(2.5,1.1){#3};\end{tikzpicture}}
\providecommand{\eg}{e.g.\ }
\providecommand{\ie}{i.e.\ }
\providecommand{\cf}{cf.\ }
\providecommand{\qbinom}[2]{\genfrac{[}{]}{0pt}{}{#1}{#2}}

\providecommand{\mathof}[1]{\mathrm{#1}}

\numberwithin{equation}{section}


\def\endappendix{}


\newcommand{\N}{\mathbb{N}}
\newcommand{\R}{\mathbb{R}}
\newcommand{\C}{\mathbb{C}}
\newcommand{\TR}{\mathbb{TR}}
\newcommand{\cod}{\mathrm{cod}}
\DeclareMathOperator{\diag}{\mathrm{diag}}
\DeclareMathOperator{\tr}{\mathrm{tr}}
\DeclareMathOperator{\SU}{SU}
\DeclareMathOperator{\GL}{GL}
\DeclareMathOperator{\SL}{SL}
\newcommand{\Rep}[1]{\mathsf{Rep}(#1)}	
\DeclareMathOperator{\Schur}{\mathsf{Schur}}
\DeclareMathOperator{\chr}{\mathsf{ch}}
\DeclareMathOperator{\WP}{\mathsf{WP}}	
\DeclareMathOperator{\interior}{int}
\DeclareMathOperator{\mult}{\mathsf{m}}

\newcommand{\newterm}[1]{\textbf{#1}}













\begin{document}
\section*{\raggedright 835: Tropical homomorphisms of the Schur-positive semiring are Newton-polytope support evaluations}
\noindent\textbf{Author:} Fritz, Tobias\par\noindent\textbf{Source:} Asymptotic and catalytic containment of representations of SU(n). Algebraic Combinatorics 7 (2024), publisher version, DOI 10.5802/alco.338. \url{https://alco.centre-mersenne.org/articles/10.5802/alco.338/}
\par\noindent\textbf{License:} CC BY 4.0, \url{https://creativecommons.org/licenses/by/4.0/}. The original publisher PDF cover has the explicit license notice. See the saved source/license records.
\par\noindent\textbf{Editorial material note (not author text):} The selected problem is Proposition3.3: every monotone tropical-valued homomorphism of the exact Schur-positive semiring with trivial kernel is a Newton-polytope support evaluation in a real direction. Full native Sections1--3 retain the Pieri/dominance proof, dominance-maximal Littlewood--Richardson additivity, concave elementary increments and complete Newton-support reconstruction, including all endpoints/kernel conventions. Named standard symmetric-function facts remain references; the neighboring positive-real classification is distinct unselected context and later quotient/catalytic-containment results are unselected. All original proof text/formulas remain editable TeX. All proof text and formulas below are editable author TeX. Mathematical wording, language and proof order are retained. Only the article class, title matter and bibliography presentation are adapted for independent compilation; no proof has been generated or repaired.
\par\noindent\textbf{Editorial difficulty:} H2: The converse proves dominance monotonicity through repeated Pieri expansions, forces additivity at the dominance-maximal Littlewood–Richardson term, and reconstructs concave elementary-function increments. These global semiring constraints yield all tropical characters; this is structurally different from positive-real evaluation and its real-root argument.
\medskip\hrule\medskip
\noindent\textbf{Author statement, complete proof and local supporting text follow in source order.}
\section{Introduction}

The goal of this note is to prove the following result:

\begin{thm}
	\label{main}
	For $n \ge 2$ and finite-dimensional complex vector spaces $V$ and $W$, let
	\[
		\rho \: : \: \SU(n) \longrightarrow \GL(V), \qquad \sigma \: : \: \SU(n) \longrightarrow \GL(W)
	\]
	be continuous representations of $\SU(n)$, or equivalently algebraic representations of $\SL(n,\C)$.
	Suppose that the following conditions hold:
	\begin{enumerate}[label=(\roman*)]
		\item\label{main_real} The associated characters $\chr$ satisfy
			\begin{equation}
				\label{main_real_ineq}
				\chr_\rho(x) \,<\, \chr_\sigma(x)
			\end{equation}
			for all $x \in \R_+^n$ with $x_1 \cdots x_n = 1$, and
		\item\label{main_trop} For the associated weight polytopes $\WP$, we have
			\begin{equation}
				\label{main_trop_ineq}
				\WP(\rho) \,\subseteq\, \interior(\WP(\sigma)),
			\end{equation}
			where $\interior$ denotes interior.
	\end{enumerate}
	Then the following subrepresentation relations hold:
	\begin{enumerate}[label=(\alph*)]
		\item \newterm{Catalytic}: there is a nonzero representation $\eta$ such that
			\begin{equation}
				\label{main_cat}
				\rho \otimes \eta \hookrightarrow \sigma \otimes \eta.
			\end{equation}
		\item \newterm{Asymptotic}: if $\sigma$ is generic, then
			\begin{equation}
				\label{main_asymp}
				\rho^{\otimes n} \hookrightarrow \sigma^{\otimes n} \qquad \forall n \gg 1.
			\end{equation}
	\end{enumerate}
	Conversely, if \eqref{main_cat} holds for some nonzero $\eta$ or \eqref{main_asymp} holds for some $n \ge 1$, then the inequalities in~\ref{main_real} hold at least non-strictly, and similarly for~\ref{main_trop} in the sense that $\WP(\rho) \subseteq \WP(\sigma)$.
\end{thm}

In this theorem and throughout the paper, all of our representations are assumed finite-dimensional complex, and as usual continuous for $\SU(n)$ and algebraic for $\SL(n,\C)$, and we will use both of these interchangeably.
We say that a representation is \newterm{generic} if it contains both the trivial and the standard representation as subrepresentations.
Finally, the \newterm{weight polytope} $\WP(\rho)$ of a representation $\rho$ is the convex hull of the weights of $\rho$ in the vector space $\mathfrak{sl}_n^*$.

\begin{rema}
	The converse part mentioned at the end of Theorem~\ref{main} is easy to prove. We include it mainly in order to indicate that our forward direction, which gives a sufficient condition for catalytic and asymptotic containment, is generically necessary as well: the only difference between the two directions is the strictness of the inequalities (and the genericity assumption on $\sigma$ in the asymptotic case).
\end{rema}

\begin{rema}
	Condition \ref{main_trop} can be thought of as extending~\ref{main_real} to infinity, in the sense that it is equivalent to
	\begin{equation}
		\label{main_trop_ineq2}
		\lim_{r \to \infty} \sqrt[r]{\chr_\rho(x^r)} \,<\, \lim_{r \to \infty} \sqrt[r]{\chr_\sigma(x^r)}
	\end{equation}
	for all $x \in \R_+^n$ with $x_1 \cdots x_n = 1$, where the map $x \mapsto x^r$ is defined componentwise.
	This follows by the fact that each side is given by maximization over the weight polytope, as in
	\[
		\log \left( \lim_{r \to \infty} \sqrt[r]{\chr_\rho(x^r)} \right) \,=\, \max_{\alpha \in \WP(\rho)} \; \sum_{i=1}^n \alpha_i \log x_i,
	\]
	which follows from the definitions upon choosing a basis of weight vectors.
\end{rema}

We prove Theorem~\ref{main} in Section~4. 
It follows as an application of a recent \newterm{Vergleichsstellensatz}\footnote{This term was introduced in~\cite{vergleichsstellensatz} by analogy with the \emph{Nullstellensatz}, which can be thought of as a separation theorem for commutative rings, and the various \emph{Positivstellens\"atze} of real algebra, which are separation theorems for ordered rings. While the Nullstellensatz is concerned with zeros (German \emph{Nullstellen}) of polynomials and a Positivstellensatz makes statements about when a polynomial or ring element belongs to the positive cone, a Vergleichsstellensatz is concerned with comparing (German \emph{vergleichen}) elements of a preordered semiring.} proven in~\cite{vergleichsstellensatz}, which is a separation theorem for preordered semirings that we recall in Section~\ref{vss_back}.
This application is facilitated by some preliminary considerations on Schur polynomials in Section~\ref{sec_schur}.
In Section~5, we present a somewhat more explicit form of Theorem~\ref{main} for $n = 2$ as Theorem~5.1.

The main open problem is as follows.

\begin{prob}
	Formulate a generalization of Theorem~\ref{main} that holds similarly for all compact Lie groups in place of $\SU(n)$.
\end{prob}

Mainly, it is not clear to us on which set of points the corresponding inequalities~\eqref{main_real_ineq} will be required to hold in the case of a general compact Lie group.
We expect that everything else can be generalized verbatim, and also that the application of our Vergleichsstellensatz will still go through, based on suitably generalized versions of Propositions~\ref{real_schur} and~\ref{trop_schur}.

\section{Background: A Vergleichsstellensatz for preordered semirings}
\label{vss_back}

Here, we present our recent Vergleichsstellensatz~\cite[Theorem~7.15]{vergleichsstellensatz} for preordered semirings after a brief recap of the relevant definitions.
Recall that a \newterm{semiring} is like a ring, except in that the existence of additive inverses is not required. 
Throughout, all of our semirings are assumed commutative.

\begin{defi}
	A \newterm{preordered semiring} is a semiring $S$ together with a preorder relation $\le$ such that both addition and multiplication are monotone,
	\[
		a \le b \quad\Longrightarrow\quad a + c \le b + c, \quad ac \le bc.
	\]
\end{defi}

Two especially important preordered semirings are the following:
\begin{itemize}
	\item The nonnegative reals $\R_+$ with the standard ordering and algebraic structure.
	\item The \newterm{tropical reals} $\TR_+ \coloneqq (\R \cup \{-\infty\}, \max, +)$ with the standard ordering.
\end{itemize}

The following definition combines the notion of power universal element introduced as~\cite[Definition~3.27]{vergleichsstellensatz} with the special case characterization given in~\cite[Lemma~3.29]{vergleichsstellensatz}.

\begin{defi}
	Let $S$ be a preordered semiring with $1 \ge 0$. Then an element $u \in S$ is \newterm{power universal} if $u \ge 1$ and for every nonzero $a \in S$ there is $k \in \N$ such that
	\[
		a \le u^k, \qquad a u^k \ge 1.
	\]
\end{defi}

\begin{rema}
	\label{pu_rem}
	\begin{enumerate}
		\item If such $S$ has a power universal element and $1 \not\le 0$, then it follows that $S$ is both zerosumfree and has no zero divisors~\cite[Remark~3.36]{vergleichsstellensatz}.
		\item By $u \ge 1$, it is enough to check the power universality inequalities on a generating subset which includes the element $2 = 1 + 1$.
			This follows from the fact that the set of $a \in S$ that satisfy the relevant conditions is closed under addition (since it contains $2$) and under multiplication (obviously).
	\end{enumerate}
\end{rema}

Our Vergleichsstellensatz then reads as follows, where we restrict the present formulation to those statements that are of interest to us here.

\begin{thm}[{\cite[Theorem~7.15]{vergleichsstellensatz}}]
	\label{vss}
	Let $S$ be a preordered semiring with $1 > 0$ and a power universal element $u$. Suppose that nonzero $a, b \in S$ satisfy the following:
	\begin{enumerate}
		\item\label{real_spec} For all monotone homomorphisms $\phi : S \to \R_+$,
			\begin{equation}
				\label{real_spec_le}
				\phi(a) < \phi(b).
			\end{equation}
		\item\label{trop_spec} For all monotone homomorphisms $\psi : S \to \TR_+$ with $\psi(u) > 0$,
			\begin{equation}
				\label{trop_spec_le}
				\psi(a) < \psi(b).
			\end{equation}
	\end{enumerate}
	Then also the following hold:
	\begin{itemize}
		\item There is nonzero $c \in S$ such that
			\begin{equation}
				\label{cat_le}
				a c \le b c.
			\end{equation}
		\item If $b$ is power universal as well, then
			\begin{equation}
				\label{asymp_le}
				a^n \le b^n \qquad \forall n \gg 1.
			\end{equation}
	\end{itemize}
	Conversely, if \eqref{cat_le} holds for some nonzero $c$ or \eqref{asymp_le} holds for some $n \ge 1$, then the inequalities \eqref{real_spec_le} and \eqref{trop_spec_le} hold non-strictly.
\end{thm}

In~\cite{vergleichsstellensatz} we also gave a concrete form for a ``catalyst'' $c$ as in~\eqref{cat_le}, but this concrete form will be of less interest to us here.
A first application of a related Vergleichsstellensatz to the asymptotics of random walks on topological abelian groups has been given in~\cite{random_walks}.
The present paper presents our second application.

\section{The preordered semiring of Schur positive polynomials}
\label{sec_schur}

Fixing a number of variables $n \in \N$, we write $\Schur_n$ for the semiring of Schur positive symmetric integer polynomials in $n$ variables, preordered such that $p \le q$ if and only if $q - p$ is Schur positive.

Equivalently, $\Schur_n$ is the preordered semiring consisting of formal sums $\sum_\lambda p_\lambda s_\lambda$ of basis elements $s_\lambda$ indexed by partitions $\lambda = (\lambda_1 \ge \cdots \ge \lambda_\ell > 0)$ of length $\ell \le n$ and with coefficients $p_\lambda \in \N$. These multiply via the unique bilinear extension of
\begin{equation}
	\label{LR_mult}
	s_\mu s_\nu \,=\, \sum_\lambda c^\lambda_{\mu\nu} s_\lambda
\end{equation}
with Littlewood-Richardson coefficients $c^\lambda_{\mu\nu}$, and the preorder is the coefficientwise one.
The unit element of $\Schur_n$ is $1 = s_{()}$, the Schur polynomial associated to the empty partition.
For $0 \le j \le n$, we also write $e_j = s_{\underbrace{\scriptstyle{(1,\ldots,1)}}_{j\:\mathrm{times}}}$ for the corresponding elementary symmetric polynomial.

\begin{lemm}
	\label{schur_bound}
	For every nonempty partition $\lambda$,
	\[
		s_\lambda \le e_1^{|\lambda|}.
	\]
\end{lemm}

\begin{proof}
	By the Pieri rule, we have $c^\lambda_{(1)\, \nu} = 1$ if and only if $\lambda$ arises from $\nu$ by adding a single box, \ie{}if there is a row index $j$ such that $j = 0$ or $\nu_{j-1} > \nu_j$, and
	\[
		\lambda_i = \begin{cases}	\nu_i + 1 & \textrm{ if } i = j, \\
						\nu_i & \textrm{ if } i \neq j. \end{cases}
	\]
	We now prove the claim by induction on the size of $\lambda$.
	The statement is trivial for $\lambda = (1)$, which is the smallest nonempty partition, since $s_{(1)} = e_1$.
	For the induction step, let $\nu$ be such that $\lambda$ arises from $\nu$ by adding a single box as above.
	Then assuming~$s_\nu \le e_1^{|\nu|}$ by induction hypothesis gives
	\[
		s_\lambda \le s_\nu e_1 \le e_1^{|\nu|+1},
	\]
	where the first inequality is by $c^\lambda_{(1)\, \nu} = 1$.
	This coincides with the claim by $|\lambda| = |\nu| + 1$.
\end{proof}

\begin{prop}
	\label{real_schur}
	The monotone homomorphisms $\Schur_n \to \R_+$ with trivial kernel can be parametrized by $y \in \R^n_{>0}$ and are given by the evaluation maps
	\begin{align*}
		\Schur_n & \longrightarrow \R_+, \\
		p & \longmapsto p(y).
	\end{align*}
\end{prop}

We thank Richard Stanley for the most difficult argument in the following proof, namely the final step of establishing $y_j \in \R$ based on~\cite{totallypositive}.\footnote{See also \url{https://mathoverflow.net/questions/419823/}.}

\begin{proof}
	The fact that these maps take nonnegative values and are monotone is immediate by the sum over semistandard Young tableaux formula for the Schur polynomials~$s_\lambda$, and it is obvious that they are homomorphisms. Also trivial kernel is obvious by $y_i > 0$ for all $i$.

	Conversely, let $\phi : \Schur_n \to \R_+$ be any monotone homomorphism with trivial kernel.
	With $\Lambda_n$ the ring of symmetric polynomials in $n$ variables, we have $\Lambda_n = \Schur_n - \Schur_n$ as the Schur polynomials form a basis.
	Therefore $\phi$ uniquely extends to a ring homomorphism $\Lambda_n \to \R$, which we also denote $\phi$ by abuse of notation, such that $\phi(s_\lambda) \ge 0$ for all partitions $\lambda$ of length $\ell(\lambda) \le n$.
	Consider now the monic single-variable real polynomial given by
	\[
		p(t) \coloneqq \sum_{i=0}^n \phi(e_i) \, t^{n-i} = \prod_{j=1}^n (t + y_j),
	\]
	where the $y_j$ are defined by the second equation, with the $-y_j$ being the roots of~$p$.
	These are a priori complex numbers where the non-real ones come in complex conjugate pairs.
	However, we claim that actually $y_j \in \R_{>0}$ for all $j$, and that $\phi$ is given by the evaluation homomorphism at $y = (y_1,\ldots,y_n)$.
	Starting with the latter, we have $\phi(e_i) = e_i(y)$ for all $i$ by Vieta's formulas, and $\phi(f) = f(y)$ for all $f \in \Lambda_n$ then follows by the fundamental theorem on symmetric polynomials.

	It remains to be shown that $y_j \in \R_{>0}$.
	By a result of Aissen and Whitney~\cite[6.]{totallypositive}\footnote{See~\cite{totallypositiveI} for the proof of this result announced in~\cite{totallypositive}.}, our claim that $p$ has real and negative roots is equivalent to the Hankel matrix of its coefficients
	\[
		H \coloneqq \left( \phi(e_{i-j}) \right)_{i,j \in \N}
	\]
	being totally positive, where our notation uses padding by zeros for the elementary symmetric polynomials with index not in $\{0,\ldots,n\}$.
	But the nonzero minors of $H$ are all of the form $\phi(s_{\lambda / \mu})$ by the second Jacobi--Trudi formula~\cite[Corollary~7.16.2]{stanley2}, and the skew Schur polynomials $s_{\lambda / \mu}$ are well-known to be Schur positive, giving $\phi(s_{\lambda / \mu}) > 0$ by the trivial kernel assumption. It follows that $H$ is indeed totally positive.
\end{proof}

Although we will not need this, it is interesting to note that $y, y' \in \R_+^n$ define the same evaluation homomorphism if and only if they differ by a coordinate permutation.
This follows by Vieta's formulas together with the fact that the coefficients of a polynomial uniquely determine its roots.

\begin{prop}
	\label{trop_schur}
	The monotone homomorphisms $\Schur_n \to \TR_+$ with trivial kernel can be parametrized by $y \in \R^n$ and are given by \smallskip
	\begin{align}
		\begin{split}
			\label{trop_schur_eq}
			\Schur_n & \longrightarrow \TR_+, \\
			\sum_{\alpha \in \N^n} p_\alpha x^\alpha & \longmapsto \max_{\alpha \in \N^n \::\: p_\alpha \neq 0} \; \sum_{i=1}^n \alpha_i y_i.
		\end{split}
	\end{align}
\end{prop}

In other words, these homomorphisms are given by linear maximization over the Newton polytope of $p$ in a fixed direction specified by $y$.
This is similar to the case of polynomial semirings~\cite[Example~2.9]{vergleichsstellensatz}.

\begin{proof}
	We first argue that every such map, denote it $\psi_y$, is a monotone homomorphism with trivial kernel.
	Monotonicity of $\psi_y$ is obvious, since $p \le q$ implies that the Newton polytope of $p$ is contained in that of $q$.
	The fact that the Newton polytope of a product of two polynomials with nonnegative coefficients is the Minkowski sum of their Newton polytopes implies that $\psi_y$ is multiplicative.
	Additivity is a consequence of the Newton polytope of a sum being the convex hull of their unions.
	Hence $\psi_y$ is a monotone homomorphism, and it has trivial kernel since it takes the value $-\infty$ only on $p = 0$.

	Conversely, suppose that $\psi : \Schur_n \to \TR_+$ is a monotone homomorphism with trivial kernel.
	We first make some preliminary observations.
	\begin{itemize}
		\item The dual Pieri rule~\cite[p.~340]{stanley2}: the product of a Schur polynomial $s_\mu$ with an elementary symmetric polynomial $e_j$ decomposes as
			\begin{equation}
				\label{se_prod}
				s_\mu e_j = \sum_\lambda s_\lambda,
			\end{equation}
			where the sum is over all partitions $\lambda$ that can be obtained from $\mu$ by adding $j$ boxes in distinct rows.
		\item Applying $\psi$ to the multiplication rule~\eqref{LR_mult} yields
			\begin{equation}
				\label{psiadd}
				\psi(s_\mu) + \psi(s_\nu) \,=\, \max_{\lambda \: : \:  c^\lambda_{\mu\nu} \neq 0} \psi(s_\lambda)
			\end{equation}
			for all partitions $\mu$ and $\nu$, where throughout we only consider partitions of length $\le n$.
		\item $\psi$ is monotone with respect to dominance ordering:
			\[
				\mu \trianglelefteq \nu \quad \Longrightarrow \quad \psi(s_\mu) \le \psi(s_\nu).
			\]
			
			To see this, we can assume without loss of generality that $\mu$ is obtained from $\nu$ by moving a single box from row $i$ down to row $j > i$, meaning that $\mu$ and $\nu$ coincide in all rows apart from
			\[
				\mu_i = \nu_i - 1, \qquad \mu_j = \nu_j + 1,
			\]
			and $j = i + 1$ (box moves down by one row and left by any number of columns) or $\mu_i = \mu_j$ (box moves left by one column and down any number of rows)~\cite[Proposition~2.3]{brylawski}.
			Then any way of adding at least $m \coloneqq \sum_{k < \nu_i} (n - \nu'_k)$ boxes to $\mu$ so as to obtain a new partition necessarily also adds the extra box at $(i, \nu_i)$.
			By repeated application of the Pieri rule, it follows that every Schur polynomial that occurs in the expansion of $s_\mu e_1^m$ also occurs in the expansion of $s_\nu e_1^m$.
			Applying $\psi$ to this statement gives
			\[
				\psi(s_\mu) + m \psi(e_1) \,\le\, \psi(s_\nu) + m \psi(e_1),
			\]
			and hence also the desired $\psi(s_\mu) \le \psi(s_\nu)$.
		\item With $\mu + \nu$ denoting the usual sum of partitions, we have\footnote{We owe this observation, which facilitates a more insightful line of proof than our original one in the next item, to an anonymous referee.}
			\[
				\psi(s_{\mu + \nu}) = \psi(s_\mu) + \psi(s_\nu).
			\]

			Indeed this follows from the previous two items upon noting that $\mu + \nu$ is dominance maximal among all partitions $\lambda$ with $c^{\lambda}_{\mu\nu} \neq 0$ by the Littlewood-Richardson rule.
		\item For every $\lambda$,
			\begin{equation}
				\label{psi_se}
				\psi(s_\lambda) \,=\, \sum_{j=1}^{\lambda_1} \psi(e_{\lambda'_j}).
			\end{equation}

			This is a direct consequence of the previous item applied to the decomposition of $\lambda$ as a sum of columns.
		\item For all $j = 1, \ldots, n$,\footnote{We again use the conventions $e_0 = 1$ and $e_{n+1} = 0$.}
			\begin{equation}
				\label{submod}
				\psi(e_{j+1}) + \psi(e_{j-1}) \le 2 \psi(e_j).
			\end{equation}
			
			Indeed by~\eqref{se_prod}, every term in the Schur polynomial expansion of $e_{j-1} e_{j+1}$ also appears in the Schur polynomial expansion of $e_j e_j$. (In fact $e_j e_j = e_{j-1} e_{j+1} + s_{\scriptstyle{(2,\ldots,2)}}$.)
	\end{itemize}
	To start the main argument, for every $f \in \Schur_n$ we have the unique decomposition into Schur polynomials
	\[
		f \,=\, \sum_\lambda f_\lambda s_\lambda
	\]
	with coefficients $f_\lambda \in \N$. Hence for the actual claim, for a given $\psi : \Schur_n \to \TR_+$ it is enough to show that there is $y \in \R^n$ such that
	\begin{equation}
		\label{psi_y}
		\psi(s_\lambda) \,=\, \psi_y(s_\lambda)
	\end{equation}
	for all $\lambda$. To this end, for every $j = 1,\ldots,n$ we put\footnote{Note that this is where the assumption that $\psi$ has trivial kernel comes in: it guarantees $\psi(e_{j-1}) > -\infty$.}
	\[
		y_j \,\coloneqq\, \psi(e_j) - \psi(e_{j-1}),
	\]
	where the second term vanishes for $j = 1$ since $e_0 = 1$.
	We then have $y_1 \ge \cdots \ge y_n$ thanks to~\eqref{submod}, and
	\begin{equation}
		\label{psiey}
		\psi(e_j) = \sum_{i=1}^j y_i
	\end{equation}
	by definition.
	By the sum over semistandard Young tableaux formula for $s_\lambda$, we can evaluate $\psi_y(s_\lambda)$ directly using the definition~\eqref{trop_schur_eq}.
	Because of $y_1 \ge \cdots \ge y_n$, the Young tableaux that attains the maximum is the minimal one, namely the one in which the $i$'th row is filled with only $i$'s for every $i = 1,\ldots,\ell(\lambda)$, so that
	\[
		\psi_y(s_\lambda) \,=\, \sum_{i=1}^{\ell(\lambda)} \lambda_i y_i.
	\]
	But then by the definition of the $y_i$ and~\eqref{psi_se}, we have
	\[
		\psi_y(s_\lambda) \,=\, \sum_{i=1}^{\ell(\lambda)} \lambda_i (\psi(e_i) - \psi(e_{i-1})) \,=\, \sum_{i=1}^{\ell(\lambda)} (\lambda_i - \lambda_{i+1}) \psi(e_i) \,=\, \sum_{j=1}^{\lambda_1} \psi(e_{\lambda'_j}) \,=\, \psi(s_\lambda),
	\]
	where the second to last equation is elementary and the last one is by~\eqref{psi_se}.
	This proves the desired~\eqref{psi_y}.
\end{proof}

\begin{rema}
	Our Theorem~\ref{vss} does not directly apply to $\Schur_n$ (or equivalently to the representation theory of $\GL(n,\C)$), since $\Schur_n$ does not have a power universal element: there is no $u$ such that $e_1 u^k \ge 1$ for any $k$, because expanding the left-hand side into Schur polynomials does not produce any constant terms.
	This is analogous to a polynomial semiring like $\R_+[x]$ not having a power universal element~\cite[Example~3.37]{vergleichsstellensatz}.
\end{rema}
\begin{thebibliography}{MMMMMM}
\bibitem[1]{totallypositive} Aissen, Michael; Edrei, Albert; Schoenberg, I. J.; Whitney, Anne On the generating functions of totally positive sequences, Proc. Nat. Acad. Sci. U.S.A., Volume 37 (1951), pp. 303-307 \href{https://alco.centre-mersenne.org/articles/10.5802/alco.338/#totallypositive}{Publisher reference}.
\bibitem[2]{totallypositiveI} Aissen, Michael; Schoenberg, I. J.; Whitney, A. M. On the generating functions of totally positive sequences. I, J. Analyse Math., Volume 2 (1952), pp. 93-103 \href{https://alco.centre-mersenne.org/articles/10.5802/alco.338/#totallypositiveI}{Publisher reference}.
\bibitem[3]{brylawski} Brylawski, Thomas The lattice of integer partitions, Discrete Math., Volume 6 (1973), pp. 201-219 \href{https://alco.centre-mersenne.org/articles/10.5802/alco.338/#brylawski}{Publisher reference}.
\bibitem[4]{random_walks} Fritz, Tobias Characterizing the asymptotic and catalytic stochastic orders on topological abelian groups, 2020 (to appear in Bernoulli) \href{https://alco.centre-mersenne.org/articles/10.5802/alco.338/#random_walks}{Publisher reference}.
\bibitem[5]{vergleichsstellensatz} Fritz, Tobias Abstract Vergleichsstellensätze for preordered semifields and semirings I, SIAM J. Appl. Algebra Geom., Volume 7 (2023) no. 2, pp. 505-547 \href{https://alco.centre-mersenne.org/articles/10.5802/alco.338/#vergleichsstellensatz}{Publisher reference}.
\bibitem[6]{stanley2} Stanley, Richard P. Enumerative combinatorics. Vol. 2, Cambridge Studies in Advanced Mathematics, 62, Cambridge University Press, Cambridge, 1999, xii+581 pages \href{https://alco.centre-mersenne.org/articles/10.5802/alco.338/#stanley2}{Publisher reference}.
\bibitem[7]{SV} Székelyhidi, László; Vajday, László Functional equations on the SU(2)-hypergroup, Math. Pannon., Volume 23 (2012) no. 2, pp. 187-193 \href{https://alco.centre-mersenne.org/articles/10.5802/alco.338/#SV}{Publisher reference}.
\end{thebibliography}
\end{document}
